% Generated draft: requires mathematical and source review.
\documentclass[11pt]{article}
\usepackage[margin=25mm]{geometry}
\usepackage{fontspec}
\setmainfont{Latin Modern Roman}
\usepackage{amsmath,amssymb,mathtools,unicode-math}
\setmathfont{Latin Modern Math}
\usepackage{enumitem,needspace,xurl,hyperref}
\setlength{\parindent}{0pt}
\setlength{\parskip}{5pt}
\setlength{\emergencystretch}{3em}
\newcommand{\N}{\mathbb{N}}
\newcommand{\Z}{\mathbb{Z}}
\newcommand{\Q}{\mathbb{Q}}
\newcommand{\R}{\mathbb{R}}
\newcommand{\C}{\mathbb{C}}
\newcommand{\F}{\mathbb{F}}
\newcommand{\A}{\mathbb{A}}
\newcommand{\PP}{\mathbb{P}}
\newcommand{\E}{\mathbb{E}}
\newcommand{\eps}{\varepsilon}
\newcommand{\dd}{\,\mathrm{d}}
\begin{document}
\section*{UnsolvedMath ID: 30002007}
\textbf{Source problem code:} catalogue entry
\subsection{Definitions and mathematical statement}
The supplied record uses \(\mathfrak h_n\) for a jet-level Cartan object and \(\mathfrak h_n^*\) for its dual, whose elements \(\theta\) are called jet-level characters. It denotes the associated Verma module by \(V(\theta)\), an associated parameter by \(\chi_\theta\), and the desired Wakimoto module by \(W(\chi_\theta)\). These are source notations: the record does not define the underlying Lie algebra, the jet construction, the allowed indices \(n\), the representation level, the module category, or the assignment \(\theta\mapsto\chi_\theta\). In particular, no specific truncated-current-algebra convention for \(\mathfrak h_n\) is assumed here.

\textbf{Statement requested by the source, conditional on specifying this framework.} For every allowed jet index \(n\) and every character \(\theta\in\mathfrak h_n^*\), construct a Wakimoto module \(W(\chi_\theta)\) together with a natural nonzero module homomorphism
\[
\Phi_\theta:V(\theta)\longrightarrow W(\chi_\theta),
\qquad \Phi_\theta\ne 0.
\]
Identify explicit additional hypotheses on \(\theta\) under which \(\Phi_\theta\) is an isomorphism. The original statement asks for suitable sufficient conditions; the cleaned statement asks to characterize the conditions for isomorphism, which may require a stronger, necessary-and-sufficient result. This difference is unresolved. Neither the additional hypotheses nor the meaning of naturality is supplied, so the displayed request is a schematic target rather than a fully specified conjecture.

The supplied dated literature assessment labels the problem partially solved and reports that the construction and a natural nonzero morphism are established for general jets, that the asserted isomorphism is known when \(n=0\), and that the higher-jet isomorphism remains conditional. The hypotheses attached to those claims are not given. These are supplied curation claims, not independently verified mathematical conclusions.
\subsection{Short English statement}
For each jet-level character \(\theta\), the source asks for a Wakimoto module and a natural nonzero map from the corresponding Verma module, together with explicit hypotheses making that map an isomorphism. The underlying representation-theoretic framework and the required hypotheses are missing. Supplied dated curation reports partial progress, including an \(n=0\) isomorphism claim whose hypotheses require review.
\subsection{Sources}
\begin{itemize}
\item Supplied local UnsolvedMath registry record, permanent ID \(30002007\), title: Verma-to-Wakimoto Isomorphisms for Jet-Level Characters. Exact source problem code: \texttt{catalogue entry}. Collection: Oberwolfach Reports Open Problems. The record gives the proposed year as 2012 and names no proposer.
\item Supplied report citation: \emph{Enveloping Algebras and Geometric Representation Theory} (2012), Oberwolfach Reports, pp.\ 733--809, DOI 10.4171/owr/2012/13. \url{https://doi.org/10.4171/owr/2012/13}. The supplied background describes these reports as documenting mathematical workshops and their open-problem discussions; no problem-specific page or speaker attribution is supplied.
\item Supplied literature reference: \emph{Compatibility of the Feigin-Frenkel and Harish-Chandra Isomorphisms}. \url{https://arxiv.org/abs/1308.2995}.
\item Supplied literature reference, labeled ``Published Compatibility Paper'': \url{https://doi.org/10.1090/S0002-9947-2014-06419-2}.
\item Literature assessment date: \texttt{2026-08-21}. This is supplied dated curation derived from the supplied OWR curation, not an independently verified assessment. It labels the problem partially solved and reports a general-jet construction and natural nonzero Verma-to-Wakimoto morphism, an isomorphism in the \(n=0\) case, and a conditional higher-jet isomorphism. The supplied background expressly requests independent verification of time-sensitive status and cited claims.
\item The registry labels the statement \texttt{corrected\_verified} and supplies a manual semantic, typographic, and LaTeX review assertion describing the cleanup as meaning-preserving, self-contained, and unsupported-content-free. Those are supplied editorial assertions; they do not resolve the missing mathematical definitions or establish mathematical correctness.
\end{itemize}
\end{document}
